% !TEX program = xelatex
\documentclass[aspectratio=169]{beamer}
\usetheme[progressbar=frametitle]{metropolis}
\usefonttheme{default}

\usepackage{graphicx}
\graphicspath{{assets/}}

\definecolor{tealblue}{RGB}{0, 128, 128}
\setbeamercolor{frametitle}{bg=tealblue, fg=white}
\setbeamercolor{progress bar}{fg=tealblue, bg=tealblue!30}
\setbeamercolor{title separator}{fg=tealblue}
\setbeamercolor{alerted text}{fg=tealblue}
\setbeamercolor{block title}{bg=tealblue, fg=white}

\title{Graphic Novel Presentation}
\subtitle{\textit{The End of the F*ing World} \& \textit{The Arrival}}
\author{[By: Andrew Ralph]}
\date{}

\begin{document}

% ── COVER ────────────────────────────────────────────────────────────────────
\maketitle


% ══════════════════════════════════════════════════════════════════════════════
% BOOK 1 — THE END OF THE F*ING WORLD
% ══════════════════════════════════════════════════════════════════════════════

% Slide 2: Title and Author
\begin{frame}{The End of the F*ing World}
  \begin{columns}[T]
    \begin{column}{0.42\textwidth}
      \includegraphics[width=\textwidth]{Cover1.png}
    \end{column}
    \begin{column}{0.55\textwidth}
      \vspace{0.5em}
      \textbf{Title:} The End of the F*ing World \\[0.6em]
      \textbf{Author:} Charles Forsman \\[0.6em]
      \textbf{Written by:} Charles Forsman \\[0.6em]
      \textbf{Illustrated by:} Charles Forsman \\[0.6em]
      \textbf{Published:} 2013
    \end{column}
  \end{columns}
\end{frame}

% Slide 3: Facts
\begin{frame}{Charles Forsman — Facts}
  \begin{columns}[T]
    \begin{column}{0.62\textwidth}
      \begin{itemize}
        \item Founded Oily Comics in 2011 --- an underground mini-comics press he ran himself
        \item His art is minimal and black-and-white. His characters are usually teenagers --- isolated, impulsive, and capable of real violence.
        \item The book was adapted into a Channel 4 series in 2017, then picked up by Netflix in 2018
      \end{itemize}
    \end{column}
    \begin{column}{0.35\textwidth}
      
\includegraphics[width=\textwidth]{CharlesForsmanCover.png}
    \end{column}
  \end{columns}
\end{frame}

% Slide 4: Other Works
\begin{frame}{Other Works}
  \begin{columns}[T]
    \begin{column}{0.55\textwidth}
      \begin{itemize}
        \item \textit{Celebrated Summer} (2013) --- Two teenagers drive around with nowhere to be
        \item \textit{Revenger} (2012) --- Someone gets wronged. They do something about it.
      \end{itemize}
    \end{column}
    \begin{column}{0.35\textwidth}
      \includegraphics[width=\textwidth]{ExtraCFCover.png}
    \end{column}
  \end{columns}
\end{frame}

% Slide 5: Genre and Audience
\begin{frame}{Genre \& Audience}
  \begin{columns}[T]
    \begin{column}{0.48\textwidth}
      \textbf{Genre}
      \begin{itemize}
        \item Dark comedy
        \item Coming-of-age
        \item Psychological thriller
      \end{itemize}
      \vspace{0.8em}
      \textbf{Intended Audience}
      \begin{itemize}
        \item Mature teens and adults (17+)
        \item Readers who connect with uncomfortable characters
      \end{itemize}
    \end{column}
    \begin{column}{0.48\textwidth}
      \textbf{Content Warnings}
      \begin{itemize}
        \item Violence/Murder and strong language
        \item Self-harm references
        \item Mature and disturbing themes
      \end{itemize}
      \vspace{0.5em}
      \includegraphics[width=\textwidth]{genre1.png}
    \end{column}
  \end{columns}
\end{frame}

% Slide 6: Story Breakdown
\begin{frame}{Story Breakdown}
  \noindent\textbf{Setting:} American suburbs. Rural roads, diners, motel rooms. \\[0.4em]
  \noindent\textbf{Characters:} James, who has decided he is a psychopath. Alyssa, who does not care what anyone thinks. \\[0.4em]
  \noindent\textbf{Inciting Incident:} James decides he wants to kill someone. He picks Alyssa. \\[0.4em]
  \noindent\textbf{Rising Action:} They run. They break into empty houses and eat whatever is left in the cabinets. James had every intention of killing her. The longer they drive, the less certain he is about that. \\[0.4em]
  \noindent\textbf{Satisfactory End:} James is shot. His last words are ``I love you.'' Alyssa carves his name into her arm.
\end{frame}

% Slide 7: Recommendation
\begin{frame}{My Recommendation}
  \begin{columns}[T]
    \begin{column}{0.38\textwidth}
      \vspace{1em}
      {\large\textbf{My Rating}} \\[0.8em]
      {\Huge\textbf{[9.7\,/\,10]}}
    \end{column}
    \begin{column}{0.58\textwidth}
      \vspace{1em}
      \textbf{I would recommend this to\ldots} \\[0.6em]
      Best read by someone who has already survived being a teenager. The alienation is specific enough to be funny in hindsight.
    \end{column}
  \end{columns}
\end{frame}


% ══════════════════════════════════════════════════════════════════════════════
% BOOK 2 — THE ARRIVAL
% ══════════════════════════════════════════════════════════════════════════════

% Slide 8: Title and Author
\begin{frame}{The Arrival}
  \begin{columns}[T]
    \begin{column}{0.42\textwidth}
      \includegraphics[width=\textwidth]{Cover2.png}
    \end{column}
    \begin{column}{0.55\textwidth}
      \vspace{0.5em}
      \textbf{Title:} The Arrival \\[0.6em]
      \textbf{Author:} Shaun Tan \\[0.6em]
      \textbf{Written by:} Shaun Tan \\[0.6em]
      \textbf{Illustrated by:} Shaun Tan \\[0.6em]
      \textbf{Published:} 2006
    \end{column}
  \end{columns}
\end{frame}

% Slide 9: Facts
\begin{frame}{Shaun Tan — Facts}
  \begin{columns}[T]
    \begin{column}{0.62\textwidth}
      \begin{itemize}
        \item Australian author and illustrator. His worlds look like nothing else on the page.
        \item Won the Academy Award for Best Animated Short Film for \textit{The Lost Thing} in 2011.
        \item Most of his work is about people who do not quite fit where they are.
      \end{itemize}
    \end{column}
    \begin{column}{0.35\textwidth}
      \includegraphics[width=\textwidth]{ShuanTuanCover.png}
    \end{column}
  \end{columns}
\end{frame}

% Slide 10: Other Works
\begin{frame}{Other Works}
  \begin{columns}[T]
    \begin{column}{0.55\textwidth}
      \begin{itemize}
        \item \textit{The Rabbits} (1998) --- animals colonize a land. It does not go well for the ones already there.
        \item \textit{Tales from Outer Suburbia} (2008) --- short stories set in neighborhoods where something is slightly off.
      \end{itemize}
    \end{column}
    \begin{column}{0.35\textwidth}
      \includegraphics[width=\textwidth]{ExtraSTCover.png}
    \end{column}
  \end{columns}
\end{frame}

% Slide 11: Genre and Audience
\begin{frame}{Genre \& Audience}
  \begin{columns}[T]
    \begin{column}{0.48\textwidth}
      \textbf{Genre}
      \begin{itemize}
        \item Wordless graphic novel
        \item Magical realism
        \item Immigration narrative
      \end{itemize}
      \vspace{0.8em}
      \textbf{Intended Audience}
      \begin{itemize}
        \item All ages --- no words to get in the way
        \item Readers drawn to history, identity, and displacement
      \end{itemize}
    \end{column}
    \begin{column}{0.48\textwidth}
      \textbf{Content Warnings}
      \begin{itemize}
        \item Themes of displacement and loss
        \item Allusions to war and persecution
      \end{itemize}
      \vspace{0.5em}
      \includegraphics[width=\textwidth]{genre2.png}
    \end{column}
  \end{columns}
\end{frame}

% Slide 12: Story Breakdown
\begin{frame}{Story Breakdown}
  \noindent\textbf{Setting:} A city with no name. Everything in it is slightly wrong in a way that makes sense. \\[0.4em]
  \noindent\textbf{Characters:} A man. His wife and daughter, left at home. \\[0.4em]
  \noindent\textbf{Inciting Incident:} He leaves. unbenknownst to the dangers ahead. \\[0.4em]
  \noindent\textbf{Rising Action:} Nothing works the way he expects. Strangers help him. Though, they have their own stories. \\[0.4em]
  \noindent\textbf{Satisfactory End:} He finds his own way. His family arrives to back him up.
\end{frame}

% Slide 13: Recommendation
\begin{frame}{My Recommendation}
  \begin{columns}[T]
    \begin{column}{0.38\textwidth}
      \vspace{1em}
      {\large\textbf{My Rating}} \\[0.8em]
      {\Huge\textbf{[9\,/\,10]}}
    \end{column}
    \begin{column}{0.58\textwidth}
      \vspace{1em}
      \textbf{I would recommend this to\ldots} \\[0.6em]
      \textit{[This is a beautifuly shown graphic novel that leaves a good bit of the story up to user interpretation. I would recommend this to anyone interested in any kind of art or dark history books]}
    \end{column}
  \end{columns}
\end{frame}

\end{document}
